\FloatBarrier


\section{Prefix Truncation}
\label{sec:truncation}

\subsection{The Hash Function}

ERE size grows with the bit-length of keys.
We need to compress keys into a smaller range $[0, 2^K)$ while preserving locality – so
that a query interval maps to a bounded number of intervals in the compressed space.
This need not be a hash in the cryptographic sense; the simplest locality-preserving
compression is plain bit-truncation:
\[
    h(x) = \lfloor x / 2^t \rfloor,
\]
keeping the top $K = L - t$ bits of each key.

Properties of the hash:
\begin{itemize}
    \item \textbf{Locality-preserving:} truncation preserves order --- if $a < b$, then
    $h(a) \leq h(b)$.
    Intervals map to intervals.
    \item \textbf{Contiguous phantoms:} recall that the phantom points of a key $x$ are
    all non-keys $y$ with $h(y) = h(x)$ --- the false pre-images that cause false positives.
    Under truncation they all fall within an interval of length~$2^t$ in the
    original universe: exactly the non-keys sharing the same $K$-bit prefix as $x$.
    This contrasts with the SODA hash, where phantoms are scattered uniformly.
\end{itemize}

\begin{figure}[h]
    \centering
    \includegraphics[width=0.85\textwidth]{phantom_comparison}
    \caption{Phantom point distribution: random hash (SODA) scatters phantoms uniformly
    across $U$; truncation concentrates them in the contiguous $2^t$-point block that contains the key.}
    \label{fig:phantom-comparison}
\end{figure}

\begin{remark}
    Throughout this section, $2^t = 2^{L - K}$ is the universe compression factor; the number of
    original-universe points that hash to the same value in $[0, 2^K)$.
    For truncation this matches the definition: $t$ bits are dropped from the back of each key.
\end{remark}

\subsection{ARE Construction With Truncation}

Combining truncation with ERE yields a filter where the contiguous phantom structure
(Figure~\ref{fig:phantom-comparison}) changes the collision analysis:

\begin{itemize}
    \item \textbf{SODA hash:} phantoms scattered $\Rightarrow$ collision zone per key is
    $\mathcal{L} \times (2^t - 1)$ (Section~\ref{sec:soda}).
    \item \textbf{Truncation:} phantoms contiguous in one block $\Rightarrow$ a query
    overlaps the block from at most $2^t - 1$ positions, giving collision zone $2^t - 1$.
\end{itemize}

The additive (rather than multiplicative) interaction eliminates the $\mathcal{L}$
factor from the filter's space formula:

\begin{table}[H]
    \centering
    \begin{tabular}{@{}llll@{}}
        \toprule
        Hash & Collision zone & $K$ for FPR $\leq \varepsilon$ & Filter BPK \\
        \midrule
        Random (SODA) & $\mathcal{L} \times (2^t-1)$ & $\log_2(n\mathcal{L}/\varepsilon)$ &
        $\log_2(\mathcal{L}/\varepsilon) + O(1)$ \\
        Truncation & $2^t - 1$ & $\log_2(2n/\varepsilon)$ &
        $\log_2(1/\varepsilon) + O(1)$ \\
        \bottomrule
    \end{tabular}
    \caption{Choosing truncation over the SODA hash reduces the resulting filter's BPK by
        $\log_2(\mathcal{L})$. At $\mathcal{L} = 128$, this is a saving of 7~bits per key.}
    \label{tab:trunc-vs-soda}
\end{table}

\begin{figure}[h]
    \centering
    \includegraphics[width=0.85\textwidth]{trunc_tradeoff_uniform}
    \caption{FPR vs BPK for truncation vs SODA on uniform keys ($n = 2^{20}$,
        $\mathcal{L} = 128$). The SODA curve is shifted right by
        ${\approx}\log_2(128) = 7$~BPK --- matching the theoretical prediction.}
    \label{fig:trunc-uniform}
\end{figure}

\FloatBarrier

\subsection{The Price of Contiguity}
\label{subsec:trunc-failure}

The same contiguity that saves space is also the weakness.

\subsubsection{FPR for Uniform Key Spacing}
\label{subsubsec:trunc-fpr}

Let $R$ denote the uniform gap between consecutive keys.
Assume keys $x_0,\, x_0{+}R,\, x_0{+}2R,\ldots$
A query $[q,\, q{+}\mathcal{L}{-}1]$ is truly empty iff it fits inside a gap:
$q = x_i + s$, $s \in (1,\, R{-}\mathcal{L}]$.

А false positive requires some key's hash value to fall inside $h([q, q+\mathcal{L}-1])$.
Since the query lies in a gap between $x_i$ and $x_{i+1}$, and $h$ is order-preserving,
all other keys are further away and their hash values fall outside this interval ---
only the two immediate neighbors can contribute.
Let $\alpha = q \bmod 2^t$ be the offset of the query start within its block.
\begin{itemize}
    \item \textbf{Left:} $x_i$ shares a block with $q$ iff $s \leq \alpha$
    (since $q = x_i + s$, both lie in the same block requires $s \leq \alpha$).
    \item \textbf{Right:} $x_{i+1} = x_i + R$ shares a block with $q{+}\mathcal{L}{-}1$ iff
    $R - s - \mathcal{L} \leq 2^t - 1 - (\alpha + \mathcal{L} - 1) \bmod 2^t$.
\end{itemize}

When $\alpha \leq 2^t - \mathcal{L}$ (query in one block), the two conditions together cover
$2^t - \mathcal{L}$ values of $s$; when $\alpha > 2^t - \mathcal{L}$ (query spans two blocks),
they cover $2{\cdot}2^t - \mathcal{L}$ values.
Averaging over $\alpha \sim \mathrm{Uniform}[0, 2^t)$:
\begin{equation}
    \Pr[\text{positive} \mid \text{empty query}]
    \;=\; \frac{(2^t - \mathcal{L})^2 + \mathcal{L}(2{\cdot}2^t - \mathcal{L})}{2^t\,(R - \mathcal{L})}
    \;=\; \frac{2^{2t}}{2^t\,(R - \mathcal{L})}
    \;=\; \frac{2^t}{R - \mathcal{L}}.
    \label{eq:trunc-fpr}
\end{equation}

Capping at 1 and handling the $R \leq \mathcal{L}$ case:
\[
    \mathrm{FPR} =
    \begin{cases}
        0 & R \leq \mathcal{L} \text{ (no empty queries)}, \\
        \min\!\bigl(1,\ 2^t/(R - \mathcal{L})\bigr) & R > \mathcal{L}.
    \end{cases}
\]

Dense or sequential keys ($R \leq \mathcal{L}$) have FPR~$= 0$ for free: every query
of length $\mathcal{L}$ hits a key, so the filter is exact.
Low FPR requires $R - \mathcal{L} \gg 2^t$; when $R - \mathcal{L} \lesssim 2^t$
the phantom block covers most of the gap and FPR approaches~1.

\subsubsection*{Breakdown thresholds for SOSD datasets}

From Eq.~\eqref{eq:trunc-fpr}, high FPR occurs when $2^t \gtrsim R - \mathcal{L}$,
so the critical truncation depth is $t^* = \lfloor\log_2(R - \mathcal{L})\rfloor$.
Table~\ref{tab:trunc-breakdown} shows $t^*$ for typical SOSD gaps and several query lengths.

\begin{table}[H]
    \centering
    \begin{tabular}{@{}llrrrr@{}}
        \toprule
        Dataset       & Typical $R$     & $\mathcal{L}=4$ & $\mathcal{L}=16$ & $\mathcal{L}=128$ & $\mathcal{L}=1024$ \\
        \midrule
        Wiki / Books  & 1--5            & $0^*$           & $0^*$            & $0^*$             & $0^*$              \\
        Facebook P50  & 28              & 4               & 3                & $0^*$             & $0^*$              \\
        Facebook mean & 387             & 8               & 8                & 8                 & $0^*$              \\
        OSM P50       & $4.5\times10^7$ & 25              & 25               & 25                & 25                 \\
        \bottomrule
        \multicolumn{6}{l}{\footnotesize ${}^{*}$~$R \leq \mathcal{L}$: every query contains a key.}
    \end{tabular}
    \caption{Critical truncation depth $t^*$ at which FPR~$\to 1$ for SOSD gap statistics
        (Table~\ref{tab:sosd-gaps}) and varying query length~$\mathcal{L}$.}
    \label{tab:trunc-breakdown}
\end{table}

Wiki and Books have gaps well below any practical $\mathcal{L}$, so every query hits a key.
Facebook's tail (gaps up to ${\sim}700$) breaks at $t \approx 8$.
OSM has $t^* \approx 25$ for all query lengths.
In a 64-bit universe this means keeping $K = 64 - 25 = 39$ bits just to avoid complete failure;
ERE compresses these to $K - \log n \approx 19$~BPK at $n = 2^{20}$.
Achieving FPR~$\leq 10^{-3}$ requires $t \leq 15$ ($K \geq 49$, $\approx 29$~BPK) under uniform queries.
Thus, truncation is impractical for clustered data.

\begin{remark}
    The formula~\eqref{eq:trunc-fpr} and the thresholds above assume \textbf{queries are placed uniformly at random in the universe} and utilize a \textbf{continuous approximation} of the discrete bit-space.
    In the exact discrete case, the numerator in Eq.~\eqref{eq:trunc-fpr} is $2^t - 1$, which is negligible for practical truncation depths.
    A more significant discrepancy arises from the query distribution: in practice, the benchmark queries of Section~\ref{sec:query-distribution} are drawn from the gaps between keys, so they land near phantom zones more often than a uniform draw would.
\end{remark}

The SODA hash avoids this entirely: pairwise independence gives
$\mathrm{FPR} \leq \varepsilon$ for any key distribution, at the cost of
$\log_2(\mathcal{L})$ extra bits per key.

\begin{figure}[h]
    \centering
    \includegraphics[width=0.85\textwidth]{trunc_tradeoff_osm}
    \caption{FPR vs BPK on OSM ($n = 2^{20}$, $\mathcal{L} = 128$, in-gap queries, Section~\ref{sec:query-distribution}).
    Truncation is bounded below by a structural FPR floor set by the gap-to-$2^t$ ratio
        (Eq.~\ref{eq:trunc-fpr}): at 24~BPK it remains at FPR~$\approx 10^{-2}$
        while SODA reaches ${\sim}10^{-5}$.}
    \label{fig:trunc-osm}
\end{figure}

\FloatBarrier

\pagebreak


\section{Adaptive Filter}
\label{sec:adaptive}

The SODA hash compresses keys into a $K$-bit universe of size $2^K = n\mathcal{L}/\varepsilon$.
If the original key spread $\max(S) - \min(S)$ is already smaller than $2^K$,
no hash is needed at all --- we build ERE directly on the (normalized) keys
and get FPR~$= 0$ for free.
The adaptive filter checks this condition at build time.

\subsection{Exact Mode}
\label{subsec:exact-mode}

\begin{enumerate}
    \item \textbf{Normalize:} subtract $\min(S)$ from all keys, shifting the dataset to
    start at~0.
    \item \textbf{Measure spread:} $M = \lceil \log_2(\max(S) - \min(S)) \rceil$.
    \item \textbf{Decide:}
    \begin{itemize}
        \item $M \leq K$: the entire dataset fits in $K$ bits $\Rightarrow$ \textbf{exact
        mode}. Build ERE directly over normalized keys. No hash, FPR~$= 0$.
        \item $M > K$: data too spread $\Rightarrow$ \textbf{SODA mode}. Apply
        a hash drawn from a pairwise-independent family, same filter guarantees as Section~\ref{sec:soda}.
    \end{itemize}
\end{enumerate}

\subsection{When Does Exact Mode Trigger?}

Exact mode requires $\max(S) - \min(S) < 2^K = n\mathcal{L}/\varepsilon$.
Rewriting in terms of density $\rho = (n - 1) / (\max(S) - \min(S))$:
\begin{equation}
    \rho > \frac{\varepsilon}{\mathcal{L}}
    \label{eq:exact-threshold}
\end{equation}
or equivalently, the average gap
$\bar{g} = (\max(S) - \min(S)) / (n - 1)$ must satisfy
$\bar{g} < \mathcal{L}/\varepsilon$.

\begin{table}[H]
    \centering
    \begin{tabular}{@{}llll@{}}
        \toprule
        $\mathcal{L}$ & $\varepsilon$ & Max avg gap        & Min density           \\
        \midrule
        128           & $10^{-3}$     & 128,000            & $7.8 \times 10^{-6}$  \\
        128           & $10^{-6}$     & $1.28 \times 10^8$ & $7.8 \times 10^{-9}$  \\
        16            & $10^{-3}$     & 16,000             & $6.25 \times 10^{-5}$ \\
        \bottomrule
    \end{tabular}
    \caption{Exact mode threshold for typical parameters. The density requirement is
    extremely low --- exact mode fails only when keys are scattered across a universe
    much wider than $n \cdot \mathcal{L} / \varepsilon$.}
    \label{tab:exact-threshold}
\end{table}

This becomes especially powerful in combination with hybrid segmentation
(Chapter~\ref{chap:hybrid}): each dense cluster is isolated into its own adaptive filter,
where the local spread is small enough for exact mode, giving FPR~$= 0$ on the dense
parts of the data.
